\documentclass[10pt,a4paper]{amsart}
\usepackage[margin=19mm]{geometry}
\usepackage{amsmath,amssymb,mathtools}
\usepackage[scr=rsfs,cal=euler]{mathalfa}
\usepackage{xfrac}
\usepackage[unicode,hidelinks,bookmarks=false]{hyperref}
\newcommand{\ie}{i.e.\ }
\newcommand{\cf}{cf.\ }
\newcommand{\resp}{respectively\ }
\newcommand{\Subsection}{\subsection}
\providecommand{\nobreakdash}{\nobreak\hbox{-}}
\setlength{\emergencystretch}{2em}
\multlinegap0pt
\usepackage{amssymb}
\usepackage{algorithm}\usepackage{algpseudocode}
\usepackage{tikz}
\usepackage{bm}
\usepackage{xcolor}
\newtheorem{prop}{Proposition}
\newtheorem{theorem}{Theorem}
\title{Shortest self-orthogonal embeddings of binary linear codes}
\author{Junmin An, Nathan Kaplan, Jon-Lark Kim, Jinquan Luo, Guodong Wang}
\date{Source: arXiv:2511.05440v2 (CC BY 4.0)}
\begin{document}
\maketitle

\noindent\emph{Source and licence.} Shortest self-orthogonal embeddings of binary linear codes. Version arXiv:2511.05440v2 (CC BY 4.0), distributed by arXiv under the Creative Commons Attribution 4.0 International licence, http://creativecommons.org/licenses/by/4.0/, as recorded in the arXiv licence metadata for this exact version and shown on the abstract page https://arxiv.org/abs/2511.05440v2. Full licence notice: \texttt{licenses/arxiv-2511.05440-CC-BY-4.0.txt}.\par\smallskip

\noindent\textit{Editorial scope.} Counted unit: the article's theorem ortho-gene (source0.tex lines 705--721). The article's complete original proof of it is delivered with it (source0.tex lines 722--763, one \texttt{proof} environment of 42 lines). No mathematics is written, repaired or re-derived: the statement and the proof are the article's own lines, in the article's own notation, and no retained line is substituted or re-flowed.  Retained uncounted as the source-local context the proof needs: lines 304--312; lines 333--339.  Nothing was omitted from any retained span, so the package carries no omission record.  No retained line contains a citation, so the wrapper carries no bibliography and no bibliography entry.  No retained line contains a figure environment, an image-inclusion command or drawing code, so no image is delivered and the package declares no asset dependency.  Editorial formatting only, in addition: an AMS \texttt{amsart} preamble that restates the article's own declarations and macros used by the retained lines, namely the environments \texttt{prop}, \texttt{theorem} on separate counters, together with the standard packages those lines need.  The retained environments print in this document as Proposition 1, Theorem 1 in order of appearance, whereas the article prints the same items with the numbering of its own full document.  The complete native source of the article as distributed in the arXiv e-print for this exact version is bundled unchanged as \texttt{sources/188665/source0.tex}.
\begin{prop} \label{thm-hull}
Let $\mathcal C$ be a binary $[n, k]$ code with generator matrix $G(\mathcal C)$ and let $\ell =\dim({\rm{Hull}}(\mathcal C))$. Then the smallest number of columns to be added to $G(\mathcal C)$ for the embedding of $\mathcal C$ into a self-orthogonal $[n', k]$ code $\tilde{\mathcal{C}}$ is
\[
\begin{cases}
 k - \ell & {\mbox{ if }} k - \ell {\mbox{ is odd, }} \\
 k - \ell {\mbox{ or }} k- \ell +1 & {\mbox{ if }} k - \ell {\mbox{ is even. }}
\end{cases}
\]
\end{prop}

\begin{theorem} \label{thm-even-odd-even}
Let $\mathcal C$ be a binary $[n, k, d]$ code which is not self-orthogonal, and let $\ell =\dim({\rm{Hull}}(\mathcal C))$.
\begin{enumerate}
    \item[(i)] If $\mathcal{C}$ is even, then the length of a shortest self-orthogonal embedding of $\mathcal{C}$ is $n+k-\ell+1$.
    \item[(ii)] If $\mathcal{C}$ is odd, then the length of a shortest self-orthogonal embedding of $\mathcal{C}$ is $n+k-\ell$.
\end{enumerate}
\end{theorem}

\begin{theorem}\label{ortho-gene}
    Let $\tilde{\mathcal{C}}$ be a shortest self-orthogonal embedding of a binary $[n, k]$ linear code $\mathcal{C}$ with $\dim(\operatorname{Hull}(\mathcal{C}))=\ell$. Let 
    \[
    \tilde{G} = \begin{bmatrix}
        G(\operatorname{Hull}(\mathcal{C})) & \mathcal{O} \\A & B
    \end{bmatrix}
    \]
    be a generator matrix for $\tilde{\mathcal{C}}$.  

    Let $\mathcal{C}'$ be the code generated by the matrix 
    \[
    G' = \begin{bmatrix}
        G(\operatorname{Hull}(\mathcal{C})) & \mathcal{O} \\A & D
    \end{bmatrix}
    \]
    where $D$ has the same dimensions as $B$.  Then $\mathcal{C}'$ is a shortest self-orthogonal embedding of $\mathcal{C}$ if and only if $D = BR$ for some orthogonal matrix $R$.
\end{theorem}

\begin{proof}
    If $\mathcal{C}$ is already self-orthogonal, then the result is immediate. For the rest of the proof we assume that $\mathcal{C}$ is not self-orthogonal.
    
    We first consider the case where $D = BR$ for an orthogonal matrix $R$. Since
    \[
    G'(G')^T=\begin{bmatrix}
        \mathcal{O} & \mathcal{O}\\\mathcal{O} & AA^T+BR(BR)^T
    \end{bmatrix}=\begin{bmatrix}
        \mathcal{O} & \mathcal{O}\\\mathcal{O} & AA^T+BB^T
    \end{bmatrix}=\mathcal{O},
    \]
    $G'$ generates a shortest self-orthogonal embedding of $\mathcal{C}$.

    Conversely, suppose that $G'$ generates a shortest self-orthogonal embedding of $\mathcal{C}$. Our goal is to show that $D = BR$ for some orthogonal matrix $R$.

    Note that $DD^T=AA^T=BB^T$. Since $\operatorname{rank}(AA^T)=k-\ell$, we have $\operatorname{rank}(B)=\operatorname{rank}(D)=k-\ell$. Let $n+m$ be the length of a shortest self-orthogonal embedding of $\mathcal{C}$. By Proposition~\ref{thm-hull}, $m=k-\ell$ or $m=k-\ell+1$.

    Suppose that $m=k-\ell$. Since both $B$ and $D$ are nonsingular, set $R=B^{-1}D$. Then we have $D=BR$ and
    \[
    RR^T=B^{-1}DD^T(B^{-1})^T=B^{-1}BB^T(B^{-1})^T=I.
    \]

    Suppose next that $m=k-\ell+1$. By Theorem~\ref{thm-even-odd-even}, $\mathcal{C}$ is an even code. Thus, every row of $A$, $B$ and $D$ has even weight. Let $\mathbf{1}$ be the all-one vector of length $m$. Define
    \[
    \hat{B}=\begin{bmatrix}
        B\\\mathbf{1}
    \end{bmatrix}\quad\mbox{and}\quad\hat{D}=\begin{bmatrix}
        D\\\mathbf{1}
    \end{bmatrix}.
    \]
    Since $k-\ell$ is even and $m=k-\ell+1$ is odd, we have $\mathbf{1}\cdot\mathbf{1}=1$. Moreover, $B\mathbf{1}^T=D\mathbf{1}^T=\mathbf{0}$. Note that both $\langle B\rangle$ and $\langle D\rangle$ are even codes. Since $\mathbf{1}\notin\langle B\rangle$ and $\mathbf{1}\notin\langle D\rangle$, both $\hat{B}$ and $\hat{D}$ are nonsingular. Set $R=\hat{B}^{-1}\hat{D}$. Since $\hat{D}=\hat{B}R$, we have $D=BR$ and
    \[
    RR^T=\hat{B}^{-1}\hat{D}\hat{D}^T(\hat{B}^{-1})^T=\hat{B}^{-1}\begin{bmatrix}
        DD^T & \mathbf{0}^T\\\mathbf{0} & 1
    \end{bmatrix}(\hat{B}^{-1})^T=
    \hat{B}^{-1}\begin{bmatrix}
        BB^T & \mathbf{0}^T\\\mathbf{0} & 1
    \end{bmatrix}(\hat{B}^{-1})^T
    =\hat{B}^{-1}\hat{B}\hat{B}^T(\hat{B}^{-1})^T=I.
    \]
    Thus, $R$ is orthogonal in both cases.
\end{proof}

\end{document}
